\section{实验设置}
\label{sec:experiments}

\subsection{数据、监督与开发范围}
使用Abilene公开X01--X24记录中的144条流\cite{abilene_data}，
以及ARI-LLM固定发布版本中的GEANT处理快照，后者为10,772行、462条流\cite{ari_release}。
GEANT尚缺少已核实的原生时间戳和物理单位映射，故窗口与距离均以槽位描述。
以零起始、左闭右开行区间表示，Abilene拟合段为$[0,33335)$，开发段为$[33384,37484)$；
GEANT对应$[0,7023)$和$[7072,8322)$。
每套数据沿用512个固定拟合窗口，GEANT拟合窗口存在重叠；
开发窗长50、步长50，分别为82和25个。

选定拟合窗的完整真值用于监督，整个拟合段用于全局统计量和相关邻流图。
20\%观测率指当前窗口的输入预算，并非历史训练始终只能接触20\%真值。
开发段用于逐轮检查点选择，也已参与研究选型；历史后续范围曾被项目使用。
所有本文准确率结果据此标为\emph{开发评价}，不宣称独立未见流量测试。

\subsection{固定观测预算与指标}
三种缺失条件均保留总计20\%的数值。
均匀条件每流保留10/50个位置；不均匀条件随机一半流保留4个、另一半保留16个。
共享缺口条件保持同样的逐流数量，在稀疏与稠密流中各取近半，
为合计一半的流设置共同连续10槽缺口；并非所有流同刻完全失观测。
未受影响流沿用不均匀条件的观测位置。
这些条件始终提供本流观测支持，不能据此宣称覆盖整流完全失观测。

对每个条件合并全部开发窗的缺失位置$\Omega$，计算
\begin{equation}
 \begin{aligned}
 \mathrm{NMAE}&=\frac{\sum_{\Omega}|\widehat x-x|}{\sum_{\Omega}|x|},\\
 \mathrm{NRMSE}&=\sqrt{\frac{\sum_{\Omega}(\widehat x-x)^2}{\sum_{\Omega}x^2}}.
 \end{aligned}
 \label{eq:metrics}
\end{equation}
三条件等权平均NMAE始终为主指标，NRMSE为辅助指标。
报告两个初始化的均值和样本标准差，并保留逐种子、逐条件数值；
相关缺失元素、三个条件和训练轮次均不充当独立训练重复。

\subsection{匹配训练与顺序筛选}
三个模型使用种子41001、41002；共有参数同种子初值相同。
AdamW采用批量8、权重衰减$10^{-4}$、梯度范数上限1，FP32且关闭TF32。
窗口顺序、训练掩码和开发掩码的种子分别为51001、61001、71001，
前两者按轮产生共同日程，训练条件近似均衡。
固定前120轮学习率$10^{-3}$，第121--160轮为$10^{-4}$，不早停、不搜索其他日程。

两个新增对照从头训练160轮；SPIN+Direct从既有第120轮的\emph{末轮}检查点接续，
完整恢复模型、AdamW和随机状态，不从已选最佳点重启。
既有组合的早期训练和续训只计同一条轨迹；本次也不把接续算作额外独立重复。
该组合的前120轮来自历史开发过程，因此共同日程不消除此前研究选型的影响。
本次降学习率诊断在新增训练前固定，160轮预算完成不等于充分优化。
保存120轮预算截面；主结果从第1--160轮中选三条件平均NMAE严格改善的最早最佳点。

按预定顺序，先完成GEANT；仅当组合在两个种子上均优于两个单路径，
且末20轮中位数方向支持这一比较，才原样扩至Abilene。
这一规则用于路线筛选，没有事后百分比门槛，也不构成显著性检验。
末20轮轨迹、支持边界及缺口内外分组仅用于解释，不替换主指标。

\subsection{参考方法与完整推理成本}
Linear为逐流可见位置线性插值，边界使用端点；无观测时用该流的拟合均值。
旧四层SPIN、旧Direct+Sync及历史Full按原有预算与初始化记录作为参考，
不参与本轮匹配路径归因，也不与新对照混合计算种子统计量。
特别地，浅层与旧四层网络的速度差异不能单独归因于Direct。

成本在同一NVIDIA H20上逐模型测量，采用各数据前8个拟合窗，
仅使用seed41001的160轮预算最佳检查点。
含SPIN的两个模型内部节点分块上限为32，Direct-only不执行该分块；
每条件、每批量先暖身5次，再计时15次。
完整推理包含输入传输、全流预测、输出回传、有限值检查、非负截断及观测回填，
不包含权重加载、拟合统计量／图准备及掩码生成。
批量1报告三条件整批中位数的平均；批量8再除以8，表示每窗摊销成本，非单请求延迟。
显存报告批量8各条件最大的PyTorch已分配量，不等同整卡占用。
Linear单独测CPU上的插值与拼接，额外输出检查在计时区间外，且不显式执行非负截断；
因此其边界与设备均不同于神经模型，不据此声称跨设备加速比。
